\subsection{The actual Hecke component in the normalized space}
Keep the original odd prime $\ell$, coefficient DVR $\mathcal O$,
arithmetic frame, rank $n\in\{2,3\}$, universal ordinary determinant
ring $B_{\rm det}$, and normalized space $\mathscr S$. Write
$A_{\mathscr S}$ for its algebra generated by the matrix entries and
$B_{\mathscr S}$ for the adic closure of the degree-zero entry monomials
and ordered local characters. The prescribed map
$\Phi:B_{\rm det}\to B_{\mathscr S}$ identifies the universal
determinant and characters after the original unit-character twist
$\kappa$ and group quotient $\pi$. These are the original equations,
including their nilpotents.\cite{support-original-transfer}

\begin{theorem}\label{thm:selmer-support-seed-transfer}
Let $J$ be the actual ordinary Hecke ideal. Retain its proved
distinguished-inertia kernel condition, the equality
$\dim(B_{\rm det}/Z)=1+n[F:\mathbb Q]$ for every minimal prime
$Z$ over $J$, and the constructed absolutely irreducible arithmetic
seed in $V(J)$ with its original symmetric polarization and local
characters. Then there are primes
$\xi\subset A_{\mathscr S}$ and $Z_{\mathscr S}\subset B_{\mathscr S}$
such that the universal representation at $\xi$ is absolutely
irreducible and
\[
 \Phi(J)\subseteq Z_{\mathscr S}=\xi\cap B_{\mathscr S},
 \qquad
 \dim(B_{\mathscr S}/Z_{\mathscr S})=1+n[F:\mathbb Q].
\]
This is the original normalized-seed condition. It does not assert
that a stable point exists in a specified characteristic-$\ell$ fibre.
\end{theorem}

\begin{proof}
Choose a minimal prime $Z$ over $J$ below the arithmetic seed, and put
$C=B_{\rm det}/Z$, $K=\overline{\operatorname{Frac}(C)}$. The actual
Hecke algebra is torsion-free over the weight ring, so $C$ is a
characteristic-zero complete local domain of the stated dimension.
Chenevier's algebraically closed field determinant theorem
(Theorem~2.12 of \emph{The $p$-adic analytic space of pseudocharacters
of a profinite group and pseudorepresentations over arbitrary rings})
realizes its generic determinant by a semisimple representation $r_0$
over $K$.\cite{support-chenevier}
A nonzero $n^2$ by $n^2$ trace-Gram determinant at the seed
comes from a nonzero element of $C$ and remains nonzero over $K$.
Thus $r_0(G)$ spans $M_n(K)$ and $r_0$ is absolutely irreducible.
No residual absolute irreducibility, or unit trace discriminant, is used.
In particular
\[
 \operatorname{tr}(r_0(g)X)=0\quad\hbox{for every }g
 \quad\Longrightarrow\quad X=0.                 \tag{ST1}
\]

Untwist by the original $\kappa$. Its trace is a continuous $C$-valued
function reducing to $\sum_i\Psi^{w_i}$, with
$w=(1,-1)$ or $(1,0,-1)$. The original inertia equations and (ST1)
kill the distinguished inertia group. At an old inertia element
$g$ with $\Psi(\pi g)=1$, the original tame condition supplies
$m$ prime to $\ell$ for which the untwisted matrix $T$ satisfies
$T^m=1$. Its characteristic polynomial $P\in C[X]$ reduces to
$(X-1)^n$. Set $Q=(X^m-1)/(X-1)$. Since $Q(1)=m$ is a unit,
$\operatorname{Res}(P,Q)$ is a unit. A Bezout identity shows
$Q(T)$ is invertible, whence $(T-1)Q(T)=0$ gives $T=1$.
This argument does not require an integral basis for $T$.

By (ST1) the untwisted kernel is the intersection of the closed sets
$t(hg)=t(h)$; polarization makes it invariant under the original
conjugation. The prescribed kernel-generation law for $\pi$ therefore
gives descent to the original normalized group. Its trace $t$ and
ordered characters $\chi_{v,i}$ remain continuous and $C$-valued.
The ordinary product matrices vanish by (ST1), applied to their
original trace equations. No local character is reordered.

The polarized trace identity gives an invertible form $Q_0$ with
$r(g)^tQ_0r(cgc^{-1})=Q_0$ and $Q_0^t=\pm Q_0$.
In rank three the alternating case is impossible. In rank two it
would imply $r(g)+r(g^*)=\operatorname{tr}(r(g))I$.
Its trace against every $r(h)$ is an identity over $C$, hence at the
actual symmetric seed. There (ST1) would give
$R+J_2R^tJ_2=\operatorname{tr}(R)I$ for every image matrix $R$.
The two off-diagonal entries are $2R_{12},2R_{21}$, forcing a
diagonal seed, a contradiction. Thus $Q_0$ is symmetric.

For the original marking element $g_0$, put
$D=r(g_0)-r(g_0^*)$. It is $Q_0$-skew and
\[
 \operatorname{tr}(D^2)=2X,
 \qquad X=t(g_0^2)-t(g_0g_0^*)\in C.
\]
If $a=\Psi(g_0)$, then $X\equiv(a-a^{-1})^2$ and
$a-a^{-1}\ne0$ in the residue field. Hensel's lemma gives a unit
$\zeta\in C$ with $\zeta^2=X$ and residue $a-a^{-1}$.
The eigenvalues of $D$ are the distinct $w_i\zeta$: in rank two
its characteristic polynomial is $X^2-\zeta^2$, and in rank three
it is $X(X^2-\zeta^2)$. Choose ordered eigenvectors and normalize
their nonzero opposite pairings. This gives the original normalized
representation $\rho$ over $K$ with
\[
 \rho(g^*)=J_n\rho(g)^tJ_n,
 \qquad \rho(g_0)-\rho(g_0^*)=\zeta\operatorname{diag}(w_i).
                                                        \tag{ST2}
\]
Only $2$ and units of $C$ enter the spectral projectors, including
when $\ell=3$. The basis itself need not be integral over $C$.

Let $M_{ij}$ be the $C$-span of all $\rho(g)_{ij}$. The marked
spectral projectors $e_i$ belong to the $C$-algebra spanned by $\rho(G)$,
so multiplication and polarization give
\[
 M_{ij}M_{jk}\subseteq M_{ik},\qquad
 M_{ii}=C,\qquad M_{ij}=M_{j^*i^*},\qquad
 M_{ij}M_{ji}\subseteq\mathfrak m_C\quad(i\ne j),       \tag{ST3}
\]
where $i^*=n+1-i$. For the last two assertions, compute
$\rho(g)_{ii}=\operatorname{tr}(e_i\rho(g))$ and
$\operatorname{tr}(e_i\rho(g)e_j\rho(h))$, then reduce the trace
identities in the original diagonal residual representation.
This proves integrality of the displayed scalar products, not
integrality of all matrix entries.

Generic irreducibility implies every $M_{ji}$ contains a nonzero
$y$. Multiplication by $y$ embeds $M_{ij}$ into $C$, making it
finite. The functions $y\rho(g)_{ij}$ are traces against fixed
elements of the image algebra and hence continuous. Artin--Rees
therefore gives continuity of $g\mapsto\rho(g)_{ij}$ in the
intrinsic $\mathfrak m_C$-adic topology of $M_{ij}$.

Form the finitely generated graded domain
\[
 D_C=C[M_{ij}z^{w_i-w_j}:i\ne j]\subset K[z,z^{-1}].
\]
Every nonempty balanced product of off-diagonal generators lies
in $\mathfrak m_C$. In rank two this follows from the opposite
weights $2,-2$ in (ST3). In rank three the weight-$1$ modules
coincide, as do the weight-$-1$ modules; call them $U,V$ and
write $P,Q$ for weights $2,-2$. The four minimal balanced products
are $UV,PQ,U^2Q,V^2P$. They lie in $\mathfrak m_C$ by (ST3),
using $U^2\subseteq P$ and $V^2\subseteq Q$.
Every balanced multiset decomposes into these pairs and triples:
remove opposite pairs first, then $1+1-2$ or $2-1-1$.

Consequently $(D_C)_0=C$ and
$\mathfrak v=\mathfrak m_C+\bigoplus_{a\ne0}(D_C)_a$ is a
maximal ideal with the original residue field. The same decomposition
gives
\[
 \mathfrak m_C^r\subseteq\mathfrak v^r\cap C
       \subseteq\mathfrak m_C^{\lceil r/3\rceil}.       \tag{ST4}
\]
The $\mathfrak v$-adic completion $\widehat D_C$ contains $D_C$
injectively by Krull intersection. Equation (ST4) makes the complete
ring $C$ a closed subring with its original topology.

Replace $\rho(g)_{ij}$ by $\rho(g)_{ij}z^{w_i-w_j}$ and leave
the local characters unchanged. Diagonal conjugation preserves
all original equations; $w_{i^*}=-w_i$ preserves the polarization.
Modulo $\mathfrak v$ this is the original diagonal residual lift.
The proved entry-module continuity makes it a continuous normalized
deformation over $\widehat D_C$. Universality supplies
$\mathscr S.R\to\widehat D_C$. Every degree-zero entry monomial
and local character maps to $C$, so their whole adic closure does too.
Call this map $\beta:B_{\mathscr S}\to C$. The matrix-entry algebra
$A_{\mathscr S}$ maps to $D_C$; evaluation at $z=1$ gives
$\xi_0:A_{\mathscr S}\to K$ with
\[
 \xi_0(\mathscr S.u)=\rho,
 \qquad \xi_0\circ\mathrm{inc}=\iota_C\circ\beta,
 \qquad \iota_C:C\hookrightarrow K.
\]

The continuous local maps $\beta\Phi$ and the quotient
$q:B_{\rm det}\to C$ have exactly the same determinant polynomial
law and ordered local characters. The original universal property
gives $\beta\Phi=q$. Thus $\beta$ is surjective. Taking
$Z_{\mathscr S}=\ker\beta$ and $\xi=\ker\xi_0$ yields
\[
 B_{\mathscr S}/Z_{\mathscr S}\simeq C=B_{\rm det}/Z,
 \qquad Z_{\mathscr S}=\xi\cap B_{\mathscr S}.
\]
The representation at $\xi$ is absolutely irreducible, and the
displayed isomorphism preserves the required dimension.
\end{proof}

\begin{theorem}\label{thm:selmer-support-starting-chart}
Under the preceding hypotheses there is a closed point
$x\in\operatorname{Spec}A_{\mathscr S}$ over the closed point of
$B_{\mathscr S}$ at which an original lower matrix entry is nonzero.
In every original normalized local chart $L$ at $x$ there is a prime
$P\subset L.R$ such that
\[
 \Phi(J)\subseteq P\cap B_{\mathscr S}=Z_{\mathscr S},\qquad
 \dim(L.R/P)=n[F:\mathbb Q]+2,
\]
and the representation at $P$ is absolutely irreducible. The chart
uses its prescribed finite extension $\mathcal O'/\mathcal O$ of
complete coefficient DVRs, with
$\mathfrak m_{\mathcal O}\mathcal O'=\mathfrak m_{\mathcal O'}$,
and the specified embedding of the finite residue field $k(x)$ into
$k(\mathcal O')$.
\end{theorem}

\begin{proof}
Before evaluation at $z=1$, the preceding proof gives a surjection
$f:A_{\mathscr S}\twoheadrightarrow D_C$: its restriction to
$B_{\mathscr S}$ is the surjection $\beta$, and its matrix-entry
images generate $D_C$ over $C$. Its prime kernel $\eta$ contracts
to $Z_{\mathscr S}$. The representation over $\operatorname{Frac}(D_C)$
is conjugate to $\rho$ after extending scalars to $K(z)$, hence is
absolutely irreducible. Since its finitely many entry coefficients are
algebraic over $\operatorname{Frac}(C)$, while a nonzero entry of
nonzero weight is transcendental over that field,
\[
 \operatorname{trdeg}_{\operatorname{Frac}(C)}
       \operatorname{Frac}(D_C)=1.                 \tag{ST5}
\]

Let $T_-$ be the ideal generated by the original lower entries, whose
degrees are $-1$ or $-2$. If $T_-\subseteq\sqrt{\mathfrak m_C D_C}$,
some $T_-^r$ is contained in $\mathfrak m_CD_C$. Every monomial of
degree at most $-2r$ contains at least $r$ negative-degree generators,
so $(D_C)_j=\mathfrak m_C(D_C)_j$ for $j\leq-2r$. Each graded
component is a finite $C$-module: Dickson's lemma gives finitely many
coordinatewise minimal exponent vectors of a fixed degree, and the
remaining factors have degree zero and lie in $C$. Nakayama therefore
makes all these components zero. This contradicts the nonzero powers
of a lower entry; such an entry exists by absolute irreducibility.
Thus the lower-entry open in
$\operatorname{Spec}(D_C/\mathfrak m_CD_C)$ is nonempty. It has a
closed point $q$ with finite residue field. Pull it back along $f$ to
obtain the required $x$.

The complete local domain $C$ is excellent and universally catenary.
We use the dimension formula and formal catenarity in their
excellent local forms.\cite{support-commutative-algebra}
The dimension formula applied at $q$, with (ST5) and finite residue
extension, gives
$\dim(D_C)_q=\dim C+1=n[F:\mathbb Q]+2$.
The original local universal property identifies $L.R$ with the
completion of $\mathcal O'\otimes_{\mathcal O}A_{\mathscr S}$
at the maximal ideal selected by its residue map. The coefficient
extension is finite etale because it is unramified and the finite
residue extension is separable. After quotienting by $\eta$, we obtain
the corresponding completion of the etale local extension of $(D_C)_q$.
It is equidimensional of dimension $n[F:\mathbb Q]+2$, and its
completion has the same property, by excellence and formal catenarity.
Every minimal prime of the completion contracts to a minimal prime
before completion by flatness, and then to zero in the domain $D_C$.
Pull one back to $P\subset L.R$. Its contraction to $A_{\mathscr S}$
is exactly $\eta$. The nonzero trace-Gram discriminant survives at
$P$, proving absolute irreducibility, and $\beta\Phi(J)=0$ gives
the exact support and contraction assertions.
\end{proof}

The trace discriminant and the fractional entry modules can depend on
the arithmetic datum; no uniform denominator estimate has been obtained.
The starting prime just constructed has characteristic zero even though
its chart's closed point has characteristic $\ell$.
The remaining step is the original global component and
characteristic-$\ell$ curve argument, followed by two-sided rigidity
and the quantitative integral Selmer budget. A characteristic-zero
seed alone cannot replace it: the matrices
$\operatorname{diag}(2,1/2)$, $I+5E_{12}$ and $I+5E_{21}$ over
$\mathbb Z_5[[t_1,\ldots,t_D]]$ span the full matrix algebra generically,
whereas their reductions generate a diagonal algebra.
